\documentclass[12pt]{article}
\usepackage[utf8]{inputenc}
\usepackage[T1]{fontenc}
\usepackage[margin=1in]{geometry}
\usepackage{setspace}
\usepackage{xcolor}
\usepackage[colorlinks=true, urlcolor=blue!60!black, linkcolor=blue!60!black]{hyperref}
\setlength{\parindent}{0em}
\setlength{\parskip}{0.45em}
\renewcommand{\baselinestretch}{1.0}
\begin{document}
\begin{center}
{\Large \textit{United Public Workers v. Mitchell}, 330 U.S. 75 (1947)\par}
\vspace{0.3em}
{\normalsize Case Brief --- One-Page Edition\par}
\end{center}

\textbf{Memory Jogger.} Hatch Act upheld; abstract pre-enforcement claims dismissed; Ninth Amendment yields to granted powers (\href{https://supreme.justia.com/cases/federal/us/330/75/}{Justia}).

\textbf{Facts.} The Hatch Act barred most federal executive employees from ``any active part in political management or in political campaigns.'' Eleven employees who wished to campaign off duty, and George P. Poole, a Mint roller facing removal for serving as a ward committeeman and party paymaster, sued.

\textbf{Procedural History.} A three-judge district court upheld the Act; the Supreme Court affirmed.

\textbf{Votes.} 4--3: Reed for Vinson, Burton, and Frankfurter; Black, Douglas (in part), and Rutledge dissented; Murphy and Jackson did not participate.

\textbf{Holding.} The eleven presented no case or controversy; Poole did, and the ban validly applied to him. Once a granted power is found, Ninth and Tenth Amendment objections fail.

\textbf{Reasoning.} Courts need ``definite rights'' and ``definite prejudicial interferences.'' The Act extended limits on ``political contributions of money'' to ``political contributions of energy,'' and treating industrial workers differently was a ``matter of detail for Congress.'' Black said the Act made employees ``mere spectators''; Douglas, that a Mint roller is as remote from policy ``as a charwoman.''

\textbf{Examples.} \emph{Same-Side:} a ban on partisan candidacy. \emph{Opposite-Side:} a ban on paid nonpartisan speeches (\textit{NTEU}). \emph{Fence-Sitter:} an anti-incumbent blog.

\textbf{Critique.} Barnett and Lash say Reed erased the Ninth Amendment; others see a right-privilege relic.

\textbf{Key Quotation.} ``The powers granted by the Constitution to the Federal Government are subtracted from the totality of sovereignty originally in the states and the people. Therefore, when objection is made that the exercise of a federal power infringes upon rights reserved by the Ninth and Tenth Amendments, the inquiry must be directed toward the granted power under which the action of the Union was taken. If granted power is found, necessarily the objection of invasion of those rights, reserved by the Ninth and Tenth Amendments, must fail.'' --- Justice Reed, 330 U.S. at 95--96.

\end{document}
